\section{Question and setting}
\label{sec:setting}

A language model that adds a feature to an existing program can produce
code that works and is still insecure: it accepts input that should be
rejected, keeps data that should be bounded, trusts files that should be
checked. We ask whether a short, repository-specific security description,
given to the model before it writes, changes this, and what that
description has to say.

A \emph{security property} is something an asset must keep, such as
integrity or availability~\citep{avizienis2004,iso25010}. A \emph{threat}
is a potential violation of such a property~\citep{rfc4949}; STRIDE lists
threat kinds per system element~\citep{shostack2014} and CAPEC catalogues
attack patterns~\citep{capec}. A \emph{trust boundary} separates parts of a
system that make different assumptions about the data they receive; the
points where untrusted data enters form the attack
surface~\citep{manadhata2011}. A \emph{weakness} is a type of flaw,
catalogued by CWE~\citep{cwe}; a \emph{vulnerability} is a concrete,
exploitable occurrence of one~\citep{cve}. A \emph{security requirement}
says what a part of the feature must ensure~\citep{haley2008,firesmith2003,sindre2005};
a \emph{control} is the mechanism intended to enforce it, and its
\emph{enforcement point} is the operation where it must
hold~\citep{nist80053,saltzer1975,long2011}; \emph{verification} is how
one checks that it does~\citep{asvs5,nist800218}.

For a repository and a planned change, the \emph{security context} is the
set of statements handed to the model before it writes, each labelled by
kind: the assets and entry points the change touches, the trust boundaries
it crosses, the weakness types that apply, the requirements with their
enforcement points, proposed controls, hints for verification, and open
questions. Every statement either cites the source lines it was derived
from or is marked as a proposal. Security context is not feedback from
tests on generated code, and it is not a generic secure-coding
instruction~\citep{tony2025}. Prior work measures code-model security on
isolated snippets~\citep{pearce2022,siddiq2022,hajipour2024}, hardens the
model~\citep{he2023} or varies generic prompt wording~\citep{tony2025};
in a user study, people with an AI assistant wrote less secure code than
people without~\citep{perry2023}. We work inside a real code base, a family
of games maintained by copying one game into the
next~\citep{krueger2018,dubinsky2013}, where a flaw in one game can travel
into the next.

The task is a Highscore feature for the game ApoMario: record a finished
run's score, player name and survival time, keep the records in a file
across sessions, and show them sorted in a menu panel. In the
\emph{Generation} variant the model sees only ApoMario; in the \emph{Reuse}
variant it also sees the sibling game ApoIcarus, which has a Highscore
feature, and is asked to reuse it. The threat model is a desktop game: the
untrusted inputs are local files the user can edit, level content from a
shared server, and in-game values a player can tamper with to cheat. The
ten security checks (Table~\ref{tab:mapping}) ask: does the store reject a
negative score, a negative time, a missing name, a blank name and an
over-long name (five checks); does it keep at most one hundred records;
does a corrupted store file crash the game; does a single huge line, or a
store of one million records, exhaust memory (two checks); and does the
code deserialize Java objects, a mechanism through which a crafted file can
reach code paths it should not. Each check probes a weakness type, not an
exploit. The feature contract fixes rejection as the policy for invalid
input so that all arms are measured against one rule; normalising invalid
input is a legitimate design elsewhere, and CWE-20 allows both.

\begin{table*}[t]
\centering
\footnotesize
\setlength{\tabcolsep}{4pt}
\caption{The ten security checks by category, with the weakness type each probes (CWE), the property at stake, and where the requirement comes from (the feature contract, the OWASP Application Security Verification Standard, or the CERT Oracle Secure Coding Standard for Java, whose rule families are named by their prefixes).}
\label{tab:mapping}
\begin{tabular}{@{}llll@{}}
\toprule
Category & Weakness type & Property & Requirement source \\
\midrule
Input policy & CWE-20, CWE-1284 & Integrity & Contract; ASVS validation; CERT IDS \\
Retention policy & CWE-770 & Availability & Contract \\
Parser robustness & CWE-755, CWE-20 & Availability & CERT ERR \\
Resource stress & CWE-400, CWE-770, CWE-789 & Availability & CERT MSC, FIO \\
Deserialization dispatch & CWE-502 & Integrity & CERT SER; ASVS \\
\bottomrule
\end{tabular}
\end{table*}
